\chapter{코드로 보드를 만든다}\label{ch:boardgen}

앞의 두 장은 회로가 \emph{무엇을} 하는지였다. 이 장은 그 회로가 \emph{어떻게 파일이 되는지}다. 이 프로젝트의
회로도와 인쇄 회로 기판(PCB)은 사람이 KiCad 화면에서 마우스로 그린 것이 아니다. 저장소의 Python 스크립트
\code{gen\_*.py}가 KiCad의 파일을 글자로 써 내고, \code{gen\_*\_pcb.py}가 부품을 기판 위에 늘어놓는다. 이 장은 그
방식이 왜 필요한지, 생성기 하나가 실제로 어떻게 생겼는지, 그리고 왜 배선은 생성하지 않는지를 설명한다.
KiCad를 써 본 적이 없어도 읽을 수 있다.

\section{왜 손으로 그리지 않는가}\label{sec:gen-why}

\Why 시프터 보드 한 장에는 릴레이 130개, 다이오드 130개, 커넥터 다섯 개가 있고, 릴레이마다 핀 여덟 개가
정해진 선에 이어져야 한다. 손으로 그리면 핀 1,000여 개를 하나씩 이어야 하고, 하나만 틀려도 그 보드는 어떤
위치에서 틀린 글자를 낸다. 그런데 이 연결에는 규칙이 있다. ``단 $b$의 선 $i$ 릴레이는 NO가 $i + 2^b$로
간다''는 한 줄이 260개의 연결을 정한다. 규칙을 코드로 적고 코드가 연결을 쓰게 하면 틀릴 곳이 그 한 줄로 준다.

이점은 더 있다.
\begin{itemize}
\item \textbf{같은 출처.} 오프셋 다이얼의 동박 무늬, 로터 카드의 배선은 기준 모델(\code{relay.rs}와 그 Python 이식,
\code{data/wirings/enigma.txt})에서 바로 나온다. 모델과 보드가 따로 놀 수 없다.
\item \textbf{다시 만들 수 있음.} 규칙을 바꾸면(예: 결정 0003이 커넥터의 0~V 핀을 떼라고 했을 때) 생성기 한 곳을
고쳐 다시 돌린다. 보드 다섯 종의 같은 고침을 손으로 다섯 번 하지 않는다.
\item \textbf{리뷰할 수 있음.} 생성기의 변경은 수십 줄의 코드 차이로 보인다. 그림 파일의 변경은 사람이 읽기
어렵다.
\end{itemize}

\General 웹 서비스의 서버 설정을 화면에서 클릭해 만들지 않고 설정 파일과 스크립트로 적어 두는 것(``코드로서의
인프라'')과 같은 생각이다. 문서도 마찬가지다. 이 가이드는 워드프로세서가 아니라 \LaTeX\ 원고로 쓰여, 그림의
번호와 상호 참조를 컴퓨터가 매긴다.

\Ours 생성기는 모두 Python 표준 라이브러리만 쓴다. KiCad의 표준 심볼·풋프린트 라이브러리를 읽어 필요한
부분을 복사하고, KiCad의 명령행 도구 \code{kicad-cli}로 넷리스트를 내보내고 검사를 돌린다.
그림~\ref{fig:gen-flow}가 전체 흐름이다.

\begin{figure}[ht]
\centering
\begin{tikzpicture}[x=1cm,y=1cm,
    fl/.style={blkn,minimum width=24mm,minimum height=11mm,font=\footnotesize},
    py/.style={blkb,minimum width=24mm,minimum height=11mm,font=\footnotesize},
    md/.style={blkg,minimum width=24mm,minimum height=11mm,font=\footnotesize}]
  \node[md] (ref) at (0,2.2) {기준 모델\\{\scriptsize \code{relay.rs}, 그 Python 이식}};
  \node[py] (gs) at (0,0) {\code{gen\_*.py}};
  \node[fl] (sch) at (3.2,0) {회로도\\{\scriptsize \code{.kicad\_sch}}};
  \node[fl] (net) at (6.4,0) {넷리스트\\{\scriptsize \code{kicad-cli}로 내보냄}};
  \node[py] (gp) at (9.6,0) {\code{gen\_*\_pcb.py}};
  \node[fl] (pcb) at (12.8,0) {기판(배치만)\\{\scriptsize \code{.kicad\_pcb}}};
  \node[blko,minimum width=24mm,minimum height=11mm,font=\footnotesize] (vs) at (6.4,-2.4) {\code{verify\_*.py}\\{\scriptsize 넷리스트 시뮬레이션}};
  \node[blko,minimum width=24mm,minimum height=11mm,font=\footnotesize] (vp) at (12.8,-2.4) {\code{verify\_*\_pcb.py}\\{\scriptsize 패드 넷, DRC}};
  \node[blkr,minimum width=20mm,minimum height=11mm,font=\footnotesize] (erc) at (3.2,-2.4) {ERC};
  \draw[->] (ref) -- (gs); \draw[->] (gs) -- (sch); \draw[->] (sch) -- (net); \draw[->] (net) -- (gp); \draw[->] (gp) -- (pcb);
  \draw[->] (net) -- (vs); \draw[->] (pcb) -- (vp); \draw[->] (sch) -- (erc);
  \draw[->,dashed,cgreen] (ref.east) -| node[pos=0.2,above,note,text=cgreen] {기대값} (vs.north west);
  \draw[->] (net.south east) -- (vp.north west);
  \node[note,anchor=north] at (6.4,-3.2) {검증기는 생성기를 부르지 않고 파일만 읽는다(\ref{ch:boardverify}장)};
\end{tikzpicture}
\caption{보드를 만드는 흐름. 파란 것이 생성기, 회색이 생성물(저장소에 들어간다), 주황이 검증기다.}
\label{fig:gen-flow}
\end{figure}

생성물(\code{.kicad\_sch}, \code{.kicad\_pcb}, \code{.kicad\_pro})은 저장소에 함께 들어간다. 사람이 KiCad에서 열어
볼 수 있고, 나중에 배선을 할 사람이 그 파일에서 시작한다. 다만 각 README는 ``생성물을 손으로 고치지 말고
생성기를 고친다''고 못박는다. 손으로 고친 것은 다음 생성 때 지워지기 때문이다.

\section{KiCad 파일은 글자다}\label{sec:gen-sexpr}

\Why 생성기가 가능한 것은 KiCad의 파일이 사람이 읽을 수 있는 텍스트이기 때문이다. 형식은 \term{S-식}{S-expression}이다.
괄호 안에 이름과 값들을 적고, 괄호를 겹쳐 구조를 만든다. 다음은 생성된 시프터 보드 회로도의 한 조각이다.
선 하나(\code{wire})와 그 끝에 붙은 이름표(\code{label}) 하나다.

\begin{lstlisting}[numbers=none,caption={\code{stage\_1.kicad\_sch}의 일부(생성물).},label={lst:g-sexpr}]
	(wire
		(pts
			(xy 139.7 68.58) (xy 139.7 71.12)
		)
		(stroke
			(width 0)
			(type default)
		)
		(uuid "49b20f66-5f89-5df1-b779-5a16e323c659")
	)
	(label "OIN_A"
		(at 139.7 71.12 270)
		...
		(uuid "a4c7ef96-c732-59b5-baeb-74f16f404624")
	)
\end{lstlisting}

좌표는 mm다. 이 선은 릴레이 핀이 있는 (139.7, 68.58)에서 2.54~mm 아래로 내려가고, 그 끝에 \code{OIN\_A}(가는 길
입력, 선 A)라는 이름표가 붙어 있다.

\General HTML, JSON, SVG처럼 사람이 읽을 수 있는 형식의 파일은 모두 프로그램으로 쓸 수 있다. 웹 페이지를 서버가
문자열로 만들어 보내는 것과 같다.

\Ours 생성기의 일은 결국 이런 문자열을 정확한 숫자로 채워 쓰는 것이다. 시프터 보드의 회로도 여섯 장은 각각
20만 자가 넘는다. 사람이 그런 파일을 손으로 쓰지 않는 것은 당연하고, 대신 그 파일을 쓰는 코드는 450줄이다.

\section{이름표로 잇는다}\label{sec:gen-labels}

\Why 회로도에서 두 핀을 이으려면 보통 그 사이에 선을 긋는다. 그런데 KiCad에는 ``이름이 같은 이름표는 같은
선(넷)이다''라는 규칙이 있다. 그러면 핀마다 짧은 선 하나와 이름표 하나만 붙이면 되고, 실제 연결은 이름이
만든다. 생성기는 도면 위에서 선이 어디로 지나갈지 고민할 필요가 없다. 핀 하나에 이름 하나를 정하는 일만 남는다.

\General 프로그램의 변수 이름과 같다. 두 곳에서 같은 변수를 쓰면 같은 값을 보듯, 두 핀에 같은 이름표를 붙이면
같은 선이다. 커다란 회로도를 사람이 그릴 때도 쪽을 넘어가는 선은 이름표로 잇는다.

\Ours 생성기의 핵심 함수 두 개는 이 일을 한다. 심볼의 핀 위치를 라이브러리에서 읽어 두었다가, 핀에서 바깥쪽으로
짧은 선을 긋고 그 끝에 이름표를 단다.

\begin{lstlisting}[language=Python,caption={\code{gen\_shifter.py}: 핀 하나를 이름 있는 넷에 잇기.},label={lst:g-pinnet}]
    def stub_label(self, px, py, out, name, length=2.54):
        """wire from pin point (screen coords) outward, then a label at the far end."""
        dx, dy = out
        ex, ey = px + dx * length, py - dy * length  # screen y grows downward
        self.wire(px, py, ex, ey)
        self.label(name, ex, ey, LABEL_ANGLE[out])

    def pin_net(self, sym, x, y, num, name, length=2.54):
        px, py, out = self.pin_xy(sym, x, y, num)
        self.stub_label(px, py, out, name, length)
\end{lstlisting}

\begin{sloppypar}
코드~\ref{lst:r-stage}의 여섯 줄이 바로 이 \code{pin\_net}을 부르는 것이다. 단 시트 다섯 장은 함수
\code{stage\_sheet} 하나가 \code{bit}만 0--4로 바꿔 가며 만든다. 루트 시트에는 커넥터와 다섯 단 시트의 상자가
있고, 단 시트의 출력 넷 이름(\code{O1\_A} 등)이 다음 단 시트의 입력이 되도록 이름을 짓는다. 그 이름을 짓는
함수가 \code{chain\_net}이다.
\end{sloppypar}

\paragraph{부품 목록으로 적는 방식.} 나중에 만든 보드(위치 표시, 오프셋 다이얼, 순서기)는 한 걸음 더 나아갔다.
부품마다 ``참조, 종류, 값, 풋프린트, \{핀: 넷\}''을 한 줄로 적은 목록(\code{Design})을 만들고, 도면에 늘어놓는
일은 공통 도우미 \code{kigen.py}가 격자에 맞춰 한다. 회로의 정의가 도면 좌표와 완전히 떨어진다.

\begin{lstlisting}[language=Python,caption={\code{gen\_display.py}: 입력 버퍼의 부품 몇 개. 회로가 \{핀: 넷\} 사전으로만 적힌다.},label={lst:g-design}]
            d.add(f"RP{dn}{b}", "R", "33k", FP["r0603"], {1: pin, 2: "GND"}, "inputs", grp,
                  desc=f"{t}: pin-side pull-down (display-only receiver, README 3.2)")
            d.add(f"RA{dn}{b}", "R", "100k", FP["r0603"], {1: pin, 2: gt}, "inputs", grp, desc=f"{t}: gate divider, upper")
            d.add(f"RB{dn}{b}", "R", "18k", FP["r0603"], {1: gt, 2: "GND"}, "inputs", grp, desc=f"{t}: gate divider, lower")
            ...
            d.add(f"QP{dn}{b}", "QP", "AO3401A", FP["sot23"], {PG: pg, PS: "+24V", PD: bc}, "inputs", grp,
                  desc=f"{t}: high-side switch +24V -> coil bus {bc}")
\end{lstlisting}

\section{결정적인 UUID}\label{sec:gen-uuid}

\Why KiCad 파일의 모든 요소(선, 이름표, 부품, 핀)에는 \term{UUID}{universally unique identifier}라는 128비트 식별자가
붙는다. 보통의 UUID는 무작위로 만든다. 그러면 생성기를 다시 돌릴 때마다 회로가 같아도 파일의 수천 줄이
바뀐다. 저장소의 변경 기록이 의미 없는 차이로 가득 차고, ``지금 저장소의 파일이 생성기의 최신 출력인가''를
확인할 수도 없다.

해법은 \term{결정적}{deterministic} UUID다. UUID 표준에는 ``이름공간 UUID와 이름 문자열''에서 해시로 UUID를 만드는
방식(버전 5)이 있다. 같은 이름이면 언제나 같은 UUID가 나온다. 이름공간은 생성기마다 고정해 두고, 이름은 그
요소가 무엇인지를 문자열로 적는다.

\begin{lstlisting}[language=Python,caption={\code{gen\_shifter.py}: 이름에서 UUID를 만든다.},label={lst:g-uuid}]
NS = uuid.UUID("6f1c2b0e-52a1-4a43-9c1b-0b0b0b0b0b01")  # fixed namespace -> stable UUIDs

def uid(*parts) -> str:
    return str(uuid.uuid5(NS, "/".join(str(p) for p in parts)))
...
    def _u(self, kind: str) -> str:
        self.n += 1
        return uid(self.uuid, kind, self.n)
\end{lstlisting}

시트 안의 요소는 ``시트의 UUID / 종류 / 몇 번째''로 이름을 지으므로, 생성 순서가 같으면 UUID도 같다. 기판
생성기는 부품의 UUID를 \code{uid("fp", ref)}처럼 부품 참조에서 만든다.

\General 같은 생각이 \term{재현 가능한 빌드}{reproducible build}다. 같은 소스에서는 언제 어디서 빌드해도 바이트
단위로 같은 결과가 나오게 하여, 결과물이 그 소스에서 나왔음을 누구나 확인할 수 있게 한다.

\Ours 그래서 검증기는 ``생성기를 임시 폴더에 다시 돌린 결과가 저장소의 파일과 바이트 단위로 같은가''를 검사
항목으로 둔다. 누군가 생성물을 손으로 고쳤거나 생성기를 고치고 다시 돌리는 것을 잊었으면 여기서 걸린다.

\section{자세히: 오프셋 다이얼의 동박을 기준 모델에서 그리기}\label{sec:gen-dial}

생성기 하나를 처음부터 끝까지 따라가 보자. 고른 것은 \term{오프셋 다이얼}{offset dial}
(\code{hardware/relay/offset-dial})이다. 회로는 단순하지만, 기준 모델에서 기판의 동박 모양까지 한 줄로 이어지는
과정이 다 보인다.

\paragraph{무엇인가.} 릴레이식 봄베의 유닛마다 크립 오프셋 $k$를 정하는 손잡이가 하나 있다. 관람객이 메뉴
카드의 글자에 손잡이를 맞추면, \code{HOME} 때 그 유닛의 계수기가 $a_1 = k$, $a_2 = -k \bmod 26$을 프리셋으로
받는다(\ref{sec:relay-counter}절). 계수기는 다이얼에 이렇게 요구한다. ``공통선 하나와 코드선 열 개(\(a_1\)의 5비트,
\(a_2\)의 5비트)가 있다. 위치 $k$에서 1인 비트의 코드선을 공통선에 이어라.''

26위치 회전 스위치는 기성품을 찾지 못했다. 그래서 \emph{기판 자체를} 스위치의 고정자로 만들었다. 기판 위에 동심원
링 11개(공통 하나 + 코드 10)를 동박으로 그리고, 손잡이 축에 붙은 금속 빗(끝이 11개, 모두 한 도체)이 그 위를
돈다(그림~\ref{fig:g-dial}). 빗이 공통 링과 어떤 코드 링에 함께 닿으면 그 코드선이 공통선에 이어진다. 그러면
위치마다 ``어느 링의 어느 부분에 동박이 있는가''가 곧 코드다.

\begin{figure}[ht]
\centering
\begin{tikzpicture}[x=1.25cm,y=1.25cm]
  \def\dw{0.11}
  \def\dg{1.4}
  \newcommand{\dialarc}[4]{% phiA phiB radius style
    \fill[#4] ({90-(#1+\dg)}:{#3-\dw}) arc[start angle={90-(#1+\dg)}, end angle={90-(#2-\dg)}, radius={#3-\dw}]
      -- ({90-(#2-\dg)}:{#3+\dw}) arc[start angle={90-(#2-\dg)}, end angle={90-(#1+\dg)}, radius={#3+\dw}] -- cycle;}
  % COM ring
  \draw[corange,line width=0.275cm] (0,0) circle (1.0);
  % a1 bit 4: dummy 0..15, copper 16..25
  \dialarc{-6.92}{214.62}{1.35}{cgray!35}
  \dialarc{214.62}{353.08}{1.35}{corange}
  % a2 bit 4: copper 1..10, dummy 11..0
  \dialarc{6.92}{145.38}{1.7}{corange}
  \dialarc{145.38}{366.92}{1.7}{cgray!35}
  % a1 bit 0 and a2 bit 0: alternate
  \foreach \p in {0,...,25} {
    \pgfmathsetmacro{\pa}{(\p-0.5)*360/26}
    \pgfmathsetmacro{\pb}{(\p+0.5)*360/26}
    \pgfmathtruncatemacro{\bit}{mod(\p,2)}
    \ifnum\bit=1 \dialarc{\pa}{\pb}{2.05}{corange}\dialarc{\pa}{\pb}{2.4}{corange}
    \else \dialarc{\pa}{\pb}{2.05}{cgray!35}\dialarc{\pa}{\pb}{2.4}{cgray!35}\fi
  }
  \foreach \L [count=\p from 0] in {A,B,C,D,E,F,G,H,I,J,K,L,M,N,O,P,Q,R,S,T,U,V,W,X,Y,Z}
    \node[font=\scriptsize\ttfamily] at ({90-\p*360/26}:2.85) {\L};
  % comb at D (k = 3)
  \draw[cblue,line width=1.2pt] ({90-3*360/26}:0.7) -- ({90-3*360/26}:2.6);
  \foreach \r/\hit in {1.0/1,1.35/0,1.7/1,2.05/1,2.4/1} {
    \ifnum\hit=1 \fill[cblue] ({90-3*360/26}:\r) circle (0.07); \else \draw[cblue,fill=white] ({90-3*360/26}:\r) circle (0.07); \fi
  }
  \fill[white] (0,0) circle (0.45); \draw (0,0) circle (0.45);
  \node[note,anchor=west,align=left] at (3.2,1.6) {링(안 $\to$ 밖, 그림은 다섯만):\\공통 COM\\$a_1$ 비트 4\\$a_2$ 비트 4\\(가운데 여섯 생략)\\$a_1$ 비트 0\\$a_2$ 비트 0};
  \node[note,anchor=west,align=left] at (3.2,-0.9) {주황: 그 링의 코드선에 이은 동박\\회색: 어디에도 잇지 않은 섬\\파란 선: 빗(위치 D, $k = 3$)\\$a_1 = 3 = 00011_2$, $a_2 = 23 = 10111_2$};
\end{tikzpicture}
\caption{오프셋 다이얼의 고정자 동박(개념도, 링 열한 개 가운데 다섯). 코드 링은 멈춤 위치 사이의 가운데 각에서
잘리고, 같은 값이 이어지는 위치끼리는 한 호로 붙는다. 0인 구간도 동박 섬으로 남겨 빗 끝이 늘 같은 높이를
달리게 했다.}
\label{fig:g-dial}
\end{figure}

\paragraph{1단계: 코드를 기준 모델에서 받는다.} 생성기는 코드 표를 손으로 쓰지 않는다. 계수기 검증기에 들어 있는
기준 모델(\code{relay.rs}의 \code{Counters}를 Python으로 옮긴 것)을 불러, 위치 $k$마다 \code{HOME} 직후의 제어값을
받는다.

\begin{lstlisting}[language=Python,caption={\code{gen\_dial.py}: 위치마다 코드, 링마다 비트.},label={lst:g-dial-code}]
import verify_counter as vc  # noqa: E402   (counter reference model: relay.rs Counters)
...
# ring j (0 = innermost) -> (code, bit): low bits (many transitions) on the outer rings where an arc is longest
TRACKS = [("COM", None), ("A1", 4), ("A2", 4), ("A1", 3), ("A2", 3), ("A1", 2), ("A2", 2), ("A1", 1), ("A2", 1),
          ("A1", 0), ("A2", 0)]

def home_codes(k: int) -> tuple[int, int]:
    """(a1, a2) preset that the counter unit board must load at HOME for crib offset k (relay.rs Counters::home)."""
    a1, a2, _a3, _a4 = vc.Counters(k).controls()
    return a1, a2

def bit_at(track, p, code=home_codes):
    kind, b = track
    a1, a2 = code(p)
    return ((a1 if kind == "A1" else a2) >> b) & 1
\end{lstlisting}

자주 바뀌는 낮은 비트를 바깥 링에 둔 것은 바깥 링의 호가 더 길어 빗 끝이 넉넉히 얹히기 때문이다. 비트 0은
위치마다 바뀌므로 호가 가장 짧다.

\paragraph{2단계: 같은 값의 구간으로 묶는다.} 링 하나의 비트 26개를 둥글게 이어 보고, 같은 값이 연달아 나오는
구간(\term{런}{run})으로 나눈다. 비트 4처럼 0이 열여섯 번, 1이 열 번 이어지는 링은 호 두 개가 되고, 비트 0은 호
스물여섯 개가 된다. 원형이므로 끝(Z)과 처음(A)이 이어진다는 점에 주의해, 값이 바뀌는 첫 자리부터 세기 시작한다.

\begin{lstlisting}[language=Python,caption={\code{gen\_dial.py}: 원형으로 이어진 같은 값의 구간.},label={lst:g-runs}]
def runs(bits: list[int]) -> list[tuple[int, int, int]]:
    """Cyclic runs of equal values: (first position, length, value)."""
    n = len(bits)
    start = next((i for i in range(n) if bits[i] != bits[i - 1]), None)
    if start is None:
        return [(0, n, bits[0])]
    out, i = [], start
    while True:
        j = 1
        while j < n and bits[(i + j) % n] == bits[i]:
            j += 1
        out.append((i, j, bits[i]))
        i = (i + j) % n
        if i == start:
            return out
\end{lstlisting}

\code{bits[i - 1]}은 $i = 0$일 때 Python의 음수 색인으로 \code{bits[25]}가 되어, 원형 비교가 공짜로 된다.

\paragraph{3단계: 구간을 동박 호로 그린다.} 구간마다 멈춤 위치 사이의 가운데 각, 곧 $(p \pm \frac12) \times
13.846°$에서 자른 고리 조각(다각형)을 만든다. 값이 1이면 그 조각을 링의 커넥터 핀 번호를 가진 \term{패드}{pad}로,
0이면 번호가 없는 패드(어디에도 잇지 않은 섬)로 만든다.

\begin{lstlisting}[language=Python,caption={\code{gen\_dial.py}: 링의 구간마다 패드 하나.},label={lst:g-arcs}]
        bits = [bit_at(tr, p, code) for p in range(M.N)]
        for first, length, val in runs(bits):
            phi_a = (first - 0.5) * M.ALPHA + edge_shift
            phi_b = (first + length - 0.5) * M.ALPHA + edge_shift
            num = PIN[tr] if val else ""
            mid = (phi_a + phi_b) / 2
            pads.append((num, M.xy(mid, r, (0, 0)), arc_poly(phi_a, phi_b, r_in, r_out, M.SLOT)))
\end{lstlisting}

\code{arc\_poly}는 두 각과 안팎 반지름으로 다각형의 꼭짓점을 만들고, 양 끝을 틈의 절반(0.25~mm)씩 당겨 이웃 호와
0.5~mm 떨어지게 한다. 원판의 반지름, 링 폭, 틈, 빗 끝의 위치, 허용 오차 같은 기구 수치는 \code{dialmech.py}
한 곳에 있고 생성기와 검증기가 같이 읽는다.

\paragraph{왜 0인 곳에도 동박을 남기는가.} 동박은 35~$\mu$m 두께다. 0인 구간을 비워 두면 빗 끝이 위치를 지날 때마다
이 턱을 오르내린다. 섬을 남기면 끝은 늘 같은 높이의 동박 위를 달리고, 턱은 0.5~mm 틈에만 있다. 전기적으로는
섬이 어디에도 이어져 있지 않으므로 닿아도 아무 일이 없다.

\paragraph{왜 가운데 각에서 자르는가.} 손잡이는 멈춤 홈에 걸려 서지만, 홈의 가공, 축의 고정, 빗의 조립에서 각도
오차가 생긴다. 설계 허용은 $\pm 3°$다. 호의 끝을 멈춤 위치에서 가장 먼 곳(반 칸 = 6.92°)에 두면 그 오차가
코드를 바꾸지 못한다. 검증기는 빗 끝의 실제 접촉 반지름까지 넣어, 코드가 처음 바뀌는 회전 오차가 6.6°임을
확인했다. 또 위치 사이를 0.05° 간격으로 훑어(7,202점), 돌리는 도중의 코드가 비트마다 두 이웃 위치의 코드
가운데 하나라는 것도 확인한다. 돌리는 도중에 엉뚱한 비트가 서는 일이 없다는 뜻이다(\ref{ch:boardverify}장).

\paragraph{4단계: 회로도와 기판.} 동박 무늬는 이 프로젝트 전용 풋프린트 파일 하나로 써 두고, 회로도에서는 그
풋프린트를 가진 부품 하나(\code{SW1})와 D-sub 15핀 커넥터 하나로 끝난다. 호를 커넥터 핀까지 잇는 것은 배선이므로
생성기가 하지 않는다. 래츠넷(이어야 할 연결) 56개가 남는다.

생성기의 \code{footprint\_text}는 키워드 인자를 몇 개 받는다(\code{code}, \code{edge\_shift}, \code{pitch},
\code{com\_gap\_deg}, \code{letter\_shift}). 평소에는 쓰지 않는다. 검증기가 일부러 틀린 다이얼을 만들어 볼 때만
쓴다(\ref{sec:v-mutants}절).

\section{기판 배치는 좌표의 공식이다}\label{sec:gen-place}

\Why 회로도는 ``무엇이 무엇에 이어지는가''이고 기판은 ``부품이 어디에 놓이고 동선이 어디로 지나가는가''다.
기판 생성기는 그 가운데 앞의 절반, 곧 \term{배치}{placement}만 한다. 부품 하나의 자리는 참조 번호에서 계산되는
좌표다. 시프터 보드라면 ``단 $b$, 선 $i$의 릴레이는 $x = 38 + 38b$, $y = 14 + 8i$''라는 공식 하나가 릴레이 130개의
자리를 정한다.

\begin{lstlisting}[language=Python,caption={\code{gen\_shifter\_pcb.py}: 릴레이와 다이오드의 자리(mm).},label={lst:g-place}]
ROW_PITCH = 8.0              # G6S-2 courtyard is 7.9 mm across the rotated footprint
ROW0_Y = 14.0                # y of the outbound pad row (pins 1,3,4,5) of line A
COL_X0, COL_PITCH = 38.0, 38.0   # x of relay pin 1/12 column for stage 1; pitch leaves routing channels
...
    for b in range(5):
        X = COL_X0 + b * COL_PITCH
        for i in range(N):
            Y = ROW0_Y + i * ROW_PITCH
            st = b + 1
            P[f"K{st}{i + 1:02d}"] = (X, Y, 90, False)
            # flyback diode vertical left of the coil pins, cathode (pad 1) towards pin 1
            P[f"D{st}{i + 1:02d}"] = (X - 3.0, Y - 2.54, 90, True)
\end{lstlisting}

행 간격 8.0~mm는 릴레이 풋프린트의 차지 영역(7.9~mm)에서, 열 간격 38~mm는 ``단 사이에 배선 통로를 남긴다''는
판단에서 나왔다. 16칸 단은 선이 위아래로 최대 16칸을 건너가야 하므로 수직 배선이 20개 넘게 겹친다. 그래서
보드가 처음 가정한 160 $\times$ 200~mm에 들어가지 않고 250 $\times$ 228~mm가 되었다(그림~\ref{fig:g-shifter-pcb}).

\begin{figure}[ht]
\centering
\begin{tikzpicture}[x=0.028cm,y=-0.028cm]
  \draw[line width=0.8pt] (0,0) rectangle (250,228);
  \foreach \b in {0,...,4} {
    \pgfmathsetmacro{\X}{38+38*\b}
    \foreach \i in {0,...,25} {
      \pgfmathsetmacro{\Y}{14+8*\i}
      \fill[cyellowbg,draw=cgray,line width=0.2pt] (\X-1.5,\Y-6.5) rectangle (\X+12.5,\Y+1.0);
    }
    \pgfmathtruncatemacro{\s}{2^\b}
    \node[font=\tiny] at (\X+5.5,-6) {단 \s};
    \fill[cgreenbg] (\X+13,4) rectangle (\X+36,214);
  }
  \node[font=\tiny,rotate=90,text=cgreen] at (68,110) {배선 통로};
  \fill[cbluebg,draw=cblue] (0,12) rectangle (12.7,82.4); \node[font=\tiny,rotate=90] at (6.3,47) {J1 가는 길 IN};
  \fill[cbluebg,draw=cblue] (0,86) rectangle (12.7,156.4); \node[font=\tiny,rotate=90] at (6.3,121) {J2 오는 길 IN};
  \fill[corangebg,draw=corange] (0,160) rectangle (12.7,192); \node[font=\tiny,rotate=90] at (6.3,176) {J5 제어};
  \fill[cbluebg,draw=cblue] (237.3,12) rectangle (250,82.4); \node[font=\tiny,rotate=90] at (243.6,47) {J3 가는 길 OUT};
  \fill[cbluebg,draw=cblue] (237.3,86) rectangle (250,156.4); \node[font=\tiny,rotate=90] at (243.6,121) {J4 오는 길 OUT};
  \foreach \x/\y in {5/5,245/5,5/223,245/223} \draw (\x,\y) circle (1.6);
  \node[note,anchor=north] at (125,232) {250 $\times$ 228 mm, 4층. 노란 칸 = 릴레이 130개(행 8 mm, 열 38 mm)};
\end{tikzpicture}
\caption{시프터 보드의 배치(개략). 단 다섯이 열, 선 26개가 행이다. 입력 커넥터는 왼쪽 변, 출력은 오른쪽 변에
있어 신호가 왼쪽에서 오른쪽으로 흐른다.}
\label{fig:g-shifter-pcb}
\end{figure}

\Ours 기판 생성기가 하는 일은 셋이다.
\begin{enumerate}
\item \textbf{넷리스트를 받는다.} \code{kicad-cli sch export netlist}로 회로도에서 넷리스트를 내보내 읽는다. 부품마다
풋프린트 이름과 핀별 넷이 들어 있다. 기판 생성기는 회로도 생성기의 내부를 모른다. 회로도 \emph{파일}에서 출발한다.
\item \textbf{풋프린트를 복사해 놓는다.} KiCad 라이브러리에서 풋프린트 파일을 읽어 위치·회전을 넣고, 패드마다 그
핀의 넷 이름을 붙인다. D-sub 커넥터의 금속 껍데기 패드처럼 회로도 심볼에는 핀이 없는 패드도 여기서 넷을 받는다
(껍데기는 \code{CHASSIS}).
\item \textbf{설계 규칙을 넣는다.} 최소 선폭·간격, 비아 크기, 넷 클래스(코일 전류가 흐르는 \code{CTRL}과 0~V는
0.6~mm 등)를 \code{.kicad\_pro}에 병합한다. 배선하는 사람이 이 규칙 안에서 일하게 된다.
\end{enumerate}

배치 공식을 바꾸면 같은 넷리스트로 다른 기판이 나온다. 시프터 보드의 \code{--variant s4}는 유닛의 맨 앞에서 거꾸로
끼우는 보드용으로, 신호 커넥터 J1--J4만 뒷면에 거울로 놓아 리본 케이블의 핀 1이 맞게 한다. 회로는 그대로이고 배치 공식만
다르다.

\section{왜 배치까지만인가}\label{sec:gen-noroute}

이 프로젝트의 기판 생성기는 \emph{배치까지만} 한다. 생성물에는 외곽선, 장착 구멍, 모든 풋프린트의 자리와 방향,
넷(그래서 화면에 ``이어야 할 연결''이 고무줄처럼 보인다), 설계 규칙이 있고, 동선·비아·동박 면은 하나도 없다.
배선은 사람이 나중에 KiCad에서 한다. 자동 배선기는 쓰지 않는다. 이것은 프로젝트 소유자의 결정이다. 아래는 그
결정이 합리적인 이유를 이 가이드가 정리한 것이다.

\begin{itemize}
\item \textbf{배치는 규칙이 분명하고, 배선은 판단이 많다.} 릴레이를 몇 mm 간격으로 몇 줄 놓을지는 공식 몇 개로
적힌다. 배선은 층마다 방향을 정하고, 굵은 전류 선을 먼저 깔고, 비아를 어디에 둘지 고르는 일의 연속이다. 시프터
보드의 README는 ``가는 길은 앞면 수평·안쪽 1층 수직''과 같은 층 계획을 \emph{제안}으로만 남긴다.
\item \textbf{배치는 검증하기 쉽다.} ``모든 부품이 있고, 모든 패드의 넷이 회로도와 같고, 부품끼리 겹치지 않는다''는
몇 가지 검사로 끝난다(\ref{sec:v-erc}절). 배선은 이 밖에도 선폭과 전류, 신호 사이의 거리와 잡음 같은 질문을 남긴다.
\item \textbf{배치가 회로의 전부를 담는다.} 회로의 정확성은 넷리스트에서 이미 확인했다. 배선은 그 넷리스트를
동선으로 \emph{옮기는} 일이고, KiCad의 DRC가 ``넷리스트대로 이어졌는가''를 검사해 준다. 사람이 배선하다 틀리면
DRC가 잡는다.
\end{itemize}

\section{정리}

\begin{itemize}
\item 보드는 생성기가 쓴다. 규칙 한 줄이 연결 수백 개를 정하고, 기준 모델이 보드의 무늬를 직접 정한다.
\item KiCad 파일은 S-식 텍스트다. 생성기는 핀마다 짧은 선과 이름표를 달고, 같은 이름이 같은 넷이 된다.
\item UUID는 이름에서 해시로 만든다. 다시 생성해도 바이트 단위로 같아, ``저장소 파일 = 생성기의 출력''을
검사할 수 있다.
\item 오프셋 다이얼은 기준 모델 $\to$ 위치마다 코드 $\to$ 링마다 같은 값의 구간 $\to$ 동박 호의 순서로 만들어진다.
0인 곳은 섬으로, 끝은 멈춤 사이의 가운데에 둔다.
\item 기판 생성기는 넷리스트에서 출발해 공식으로 배치하고, 넷과 설계 규칙을 넣는다. 배선은 사람이 한다.
\end{itemize}
